\begin{abstract}
Few-shot industrial anomaly detection increasingly relies on frozen vision foundation features, especially DINOv2 patch embeddings with memory-bank or PCA/subspace scoring. These methods can rank anomalies well, but their raw distances are not reliable probabilities, and calibration becomes brittle under low-shot, class, dataset, and corruption shift. This paper studies a decoupled reliability layer for low-storage DINOv2 subspace anomaly detection. The raw anomaly ranking is kept as the PCA/subspace residual, while probability-like reliability is produced by post-hoc calibrators and leave-one-image-out (LOIO) conformal p-value views derived only from normal support images. On MVTec AD, the resulting detector is competitive but not state-of-the-art in image AUROC, while using substantially less reference storage than memory-bank baselines. On VisA, increasing the subspace dimension from 64 to 128 improves AUROC by 0.011--0.017 with storage remaining below 0.6 MB. Most importantly, on a full VisA corruption benchmark with 12 classes, $k\in\{4,8\}$, five seeds, and four corruptions, LOIO conformal reliability achieves an overall ECE of 0.0766 and improves over Vector Platt and Shift-Aware Platt in every tested corruption cell. We position the method as a calibrated, low-storage reliability-routing framework rather than an MVTec AUROC SOTA or adversarial-robustness claim.
\end{abstract}
